\section{Discussion}
\label{sec:discussion}

\paragraph{What the two results share, and where they differ.}
In both settings the agents wrote down a consequence that the terms in front of them did not show, and then acted in line with it. The debt result can be checked in a way the bargaining result cannot. Whether a debt needs unlisted mortgages follows from four sums and the rules, so the hard cases are defined without reference to what any agent did or said. In those cases the stated arithmetic was right in all but one instance, and the action matched the statement every time. The systems that left out released collateral once counted it in other games, so the two failures are lapses and do not mark a limit of any system.

In bargaining, the consequence is a judgment about value: how much a rival's monopoly will cost the holder later. The agents stated that judgment and refused accordingly, but their refusals tracked the event of completion and the form of the offer. They did not track the amount of building the rival could afford, and some stated amounts were wrong. Two systems declined almost every such opening under an announced rule or commitment. A rule of that kind may be a sound response to an opportunity that is hard to price. It is also what an agent would do if the word ``monopoly'' alone drove the refusal. The records cannot separate the two.

\paragraph{Why form might matter.}
A counter to a cash bid changes one number, while a swap requires valuing two bundles and the rival's gain together. The swaps on offer may also have been worse in ways list prices miss, and proposer mix differs across the cells. A test would hold the recipient's side of the offer fixed and vary only whether the proposer completes a group, and separately vary how much the proposer could then build.

\paragraph{The systems are configured systems.}
The response conditions are part of what was measured. Opus played without private notes after the first games, and Grok's offers include many that its own text disowned and that other agents accepted. With several agents in one game, a change to one participant's interface changes the offers every other participant sees.

\section{Related work}
\label{sec:related}

\paragraph{Negotiation with language models.}
Early neural negotiation agents were trained on dialogues with hidden rewards~\cite{lewis2017deal}. NegotiationArena evaluates LLM bargaining in defined scenarios and reports that tactics such as feigned desperation change payoffs~\cite{bianchi2024negotiationarena}. Cosentino et al.\ find that agents can model a counterpart's preferences accurately and still fail to turn that into better outcomes~\cite{cosentino2026counterparty}. Our levels of evidence follow that line. Those studies have an explicit payoff for each agreement and controlled variation, which ours lacks. Ours has persistent consequences: a trade changes who can build, who goes bankrupt and who inherits.

\paragraph{Externalities and trade restrictions.}
Economic theory treats sales in which the identity of the buyer changes the position of others. Jehiel and Moldovanu show that such externalities can delay agreement~\cite{jehiel1995delay}, and Jehiel, Moldovanu and Stacchetti study sales that affect the later interaction among potential buyers~\cite{jehiel1996nuclear}. A completing deed is an object of this kind, and the G13 sale is priced on a harm to a third party. We use the idea descriptively and test no equilibrium prediction. Refusing trades that help a rival is also not new to LLM agents. Guhe and Lascarides study persuasion and trade embargoes among symbolic agents in Settlers of Catan and measure their effect on winning in controlled simulations~\cite{guhe2014persuasion}. Our records show free-form versions of such refusals without the comparison that would say whether they help.

\paragraph{Monopoly and economic agent benchmarks.}
AgentMonopoly varies model family, size and persona in four-player Monopoly and reports trade frequency and acceptance~\cite{wu2026agentmonopoly}. FinBoardBench covers financial management across board games and finds that models buy assets while short of cash~\cite{hu2026finboardbench}. Cattle Trade pits LLMs against scripted agents and relates resource discipline to rank~\cite{muller2026cattletrade}. Li et al.\ check inter-agent messages in long commerce simulations against simulator state, with a reference set and several judge models~\cite{li2026arena}. Playing Monopoly, countering offers and managing cash are therefore not our contribution. We add a complete census of offers and debts paired with the text recorded before each action, and the split of debts by what the first menu could raise.

\paragraph{Planning and stated reasons.}
PlanBench evaluates plan generation and reasoning about change in planning domains and reports that LLMs fall well short on plan generation~\cite{valmeekam2023planbench}. Debt payment here is a different task. Execution has feedback because the menu refreshes after each action. The decision to continue has none, and the sixteen cases place the test there. Explanations produced by language models need not reflect the process that produced the answer~\cite{turpin2023unfaithful,lanham2023faithfulness}, so we compare each note with the state and the action and do not treat it as evidence of how the action was chosen.

\section{Limitations}
\label{sec:limits}

The cohort is twenty games from five seeds, one roster and one cyclic seat order. It cannot separate a system from its seat, its opponents or its luck, and it supports no ranking. Wins were \fact{p.wins-by-system.Sol}, \fact{p.wins-by-system.Gemini}, \fact{p.wins-by-system.Opus} and \fact{p.wins-by-system.Grok} for Sol, Gemini, Opus and Grok, and \fact{p.wins-by-seat.1}, \fact{p.wins-by-seat.2}, \fact{p.wins-by-seat.3} and \fact{p.wins-by-seat.4} by seat.

Every contrast was selected after the games were inspected, and no new games were run to confirm any of them. The sixteen hard debts include only two failures. The offer comparisons differ in terms, counterparties and timing as well as in completion and form. No reference policy, value estimate or human annotation is available, so nothing here says whether a refusal, a price or a liquidation was good play. Whether a rule that takes any listed financing action would pay the same debts is untested. Such a rule would never face the decision to quit, where both failures occurred.

The readings of text come from language-model assistants, some from the same model families as two participants, and exact spans are printed so that a reader can check them. The three payment sequences behind the payable bankruptcies and the replay comparisons are inherited from earlier work in this project. The prompt invites public deception, so the incidents show permitted conduct. The response conditions for Opus and Grok shaped what every player faced.

\section{Conclusion}
\label{sec:conclusion}

Across twenty complete games, LLM agents stated what a move would make possible before they made it, and their actions agreed with those statements. For their own debts the statements were accurate. Agents almost always quit when no payment existed and continued when one did. They paid fourteen of the sixteen debts that could be met only with mortgages absent from the menu, and in thirteen of those they had counted that collateral before acting. For rivals the statements were as frequent and less exact. Agents almost never accepted an opening that completed only the proposer's group, they refused swaps and put prices on cash bids, and their refusals followed the event of completion without following how much the rival could then build. Whether these choices were good ones, and whether they hold in a new cohort, are questions for controlled tests.
